\documentclass[11pt]{article}

\usepackage[T1]{fontenc}
\usepackage{lmodern}
\usepackage{amsmath,amssymb,amsthm,mathtools}
\usepackage{microtype}
\usepackage[margin=0.92in]{geometry}
\usepackage{booktabs,longtable,array}
\usepackage{enumitem}
\usepackage{xcolor}
\usepackage{listings}
\usepackage{hyperref}
\usepackage{cleveref}
\usepackage{fancyhdr}
\usepackage{lastpage}

\definecolor{BusyBlue}{RGB}{28,67,118}
\definecolor{BusyGray}{RGB}{90,96,104}
\definecolor{CodeBg}{RGB}{248,249,251}
\definecolor{CodeRule}{RGB}{205,211,220}

\hypersetup{
  colorlinks=true,
  linkcolor=BusyBlue,
  urlcolor=BusyBlue,
  citecolor=BusyBlue,
  pdftitle={Contender 14: Self-Height Operator Power over Dynamic-Kind System F-omega},
  pdfauthor={Research package for the BusyLean project}
}

\pagestyle{fancy}
\fancyhf{}
\lhead{BusyLean Contender 14}
\rhead{Research proposal}
\cfoot{\thepage\ / \pageref*{LastPage}}
\setlength{\headheight}{14pt}
\setlength{\emergencystretch}{2em}

\lstdefinelanguage{lean}{
  keywords={def,theorem,inductive,structure,where,match,with,fun,let,in,
    namespace,end,open,abbrev,import,by,have,show,exact,termination_by,
    cases,unfold,rw,change,intro},
  sensitive=true,
  morecomment=[l]{--},
  morecomment=[s]{/-}{-/},
  morestring=[b]"
}
\lstset{
  language=lean,
  basicstyle=\ttfamily\footnotesize,
  keywordstyle=\bfseries\color{BusyBlue},
  commentstyle=\itshape\color{BusyGray},
  backgroundcolor=\color{CodeBg},
  frame=single,
  framerule=0.35pt,
  rulecolor=\color{CodeRule},
  columns=fullflexible,
  keepspaces=true,
  showstringspaces=false,
  breaklines=true,
  breakatwhitespace=false,
  xleftmargin=0.6em,
  xrightmargin=0.6em,
  aboveskip=0.8em,
  belowskip=0.8em,
  literate=
    {→}{{$\to$}}1 {←}{{$\leftarrow$}}1 {ℕ}{{$\mathbb{N}$}}1
    {∀}{{$\forall$}}1 {∃}{{$\exists$}}1 {∈}{{$\in$}}1 {≤}{{$\le$}}1
    {≥}{{$\ge$}}1 {×}{{$\times$}}1 {·}{{$\cdot$}}1
    {⟨}{{$\langle$}}1 {⟩}{{$\rangle$}}1 {ω}{{$\omega$}}1
}

\newtheorem{theorem}{Theorem}[section]
\newtheorem{lemma}[theorem]{Lemma}
\newtheorem{proposition}[theorem]{Proposition}
\newtheorem{corollary}[theorem]{Corollary}
\theoremstyle{definition}
\newtheorem{definition}[theorem]{Definition}
\newtheorem{construction}[theorem]{Construction}
\theoremstyle{remark}
\newtheorem{remark}[theorem]{Remark}
\newtheorem{warning}[theorem]{Validation note}

\newcommand{\N}{\mathbb{N}}
\newcommand{\Fd}{\mathsf{Fd}}
\newcommand{\Fam}{\mathsf{Family}}
\newcommand{\Prog}{\mathsf{Progressive}}
\newcommand{\OpStage}{\mathsf{operatorStage}}
\newcommand{\OpPow}{\mathsf{operatorPower}}
\newcommand{\lean}[1]{\begingroup\def\_{\textunderscore\allowbreak}\ifmmode\text{\ttfamily #1}\else\texttt{#1}\fi\endgroup}
\newcommand{\Cthirteen}{C_{13}}
\newcommand{\Cfourteen}{C_{14}}
\newcommand{\Fthirteen}{F_{13}}
\newcommand{\Ffourteen}{F_{14}}

\title{\textbf{Contender 14: Self-Height Operator Power\\
  over Dynamic-Kind System F$\boldsymbol{\omega}$}\\[0.55em]
  \large A constructive large-number proposal for BusyLean}
\author{Research package for the BusyLean project}
\date{October 2026}

\begin{document}
\maketitle

\begin{abstract}
The current BusyLean incumbent, Contender~13, takes one bounded reflection step
through a strongly normalizing, dynamic-kind extension of System
F$\omega$.  We propose a successor that changes the \emph{amplifier}, not the
underlying normalization proof.  First, the unary Contender~13 reflection is
lifted to countable oracle families while retaining explicit domination of
every visible family coordinate.  Second, this progressive transformer is
closed under BusyLean's operator-stage construction along a schedule obtained
by placing another autonomous critical-point closure above Contender~12's
schedule.  Third, the complete operator stage is iterated a self-fed finite
number of times: at input $n$, the height is the previous full family value
$G(n)+1$, where $G(n)=F_{13}(F_{13}(n))$.  Finally,
\[
  F_{14}(n)=P_{G(n)+1}(\underline{G})(G(n),G(n)),
  \qquad C_{14}=F_{14}(F_{14}(42)).
\]
Every operation is transparent, structural, and executable in principle.  The
supplied Lean development proves the pointwise margin $G(n)<F_{14}(n)$, the
required comparison $C_{13}<C_{14}$, and domination of every accepted
Fd-program at the construction's diagonal budget.  The article separates
these formal claims from heuristic proof-theoretic calibration and records the
exact validation status of the package.
\end{abstract}

\vspace{0.35em}
\noindent\fbox{\parbox{0.965\textwidth}{\raggedright
\textbf{Main proposed theorem.}
\(
  \lean{BusyLean.contender\_13\_lt\_contender\_14} :
  \lean{contender\_13 < contender\_14}.
\)
The proof is reduced to explicit occurrence inequalities in already established
reflection machinery; it does not compare or evaluate either gigantic numeral.
}}

\tableofcontents
\newpage

%======================================================================
\section{Repository snapshot and design target}

This work targets BusyLean at commit
\texttt{0cd13c0c2607b4fcb461090a95f356e538c0fc5b}, whose documented incumbent is
Contender~13~\cite{busyLean}.  The repository is a Lean~4 version of the
constructive large-number contest initiated in Coq by Cody Roux and inspired by
Aaronson's larger-number essay~\cite{rouxContest,aaronson1999}.  Its relevant
acceptance conditions are unusually stringent:

\begin{enumerate}[leftmargin=*,label=\arabic*.]
\item the number is a closed, computable \lean{Nat};
\item the strict comparison with the previous contender is a formal theorem;
\item no \lean{noncomputable}, \lean{partial}, \lean{unsafe}, opaque stand-in,
      external implementation, \lean{native\_decide}, or custom axiom is used;
\item only the repository's permitted logical axioms may occur in the proof
      closure;
\item the definition may use the previous contender's \emph{construction} as
      a seed or oracle, but not merely mention the previous closed constant and
      add one; and
\item the definition and proof must remain within the repository's checking
      budgets.
\end{enumerate}

The pinned toolchain is Lean~4.34.0.  Lean is both the implementation language
and the trusted proof checker~\cite{deMouraUllrich2021}.  Consequently, the
central design problem is not to write down a huge informal expression.  It is
to expose a transparent total program whose strict advantage is witnessed by a
small proof term that never asks the elaborator to normalize the number.

\subsection{The incumbent in one equation}

Write $R^{\Fd}_f(n)$ for the repository function \lean{Fd.reflectFd f n}.  It
is definitionally
\begin{equation}\label{eq:fd-reflection}
  R^{\Fd}_f(n)
  = n+1+\max\left(f(n),
    \max_{t\in \mathcal{T}_{\Fd}(n)} \operatorname{value}_f(n,t)\right),
\end{equation}
where $\mathcal{T}_{\Fd}(n)$ is the finite list of dynamic-kind terms at budget
$n$.  The total evaluator is justified by the repository's strong-normalization
proof for the object language.  Contender~13 is
\begin{equation}\label{eq:c13}
  F_{13}(n)=R^{\Fd}_{\mathrm{family}_{12}}(n),
  \qquad C_{13}=F_{13}(F_{13}(42)).
\end{equation}
The language includes a predicative dynamic kind that packages operators of
all permitted dyn-free kinds.  It is therefore materially stronger than the
kind-polymorphic F$\omega$ language used by Contenders 10--12.  Any credible
successor should exploit this stronger language rather than retreat to an
older base.

%======================================================================
\section{Rejected shortcuts}

The following constructions would establish a numerical inequality but do not
meet the research objective.

\paragraph{Adding one.}
Defining $C_{14}=C_{13}+1$ is constructive and trivially larger, but violates
the no-laziness rule and contributes no structural growth mechanism.

\paragraph{A bare extra diagonal.}
Replacing $F_{13}(F_{13}(42))$ by a few more literal nestings of $F_{13}$ is
better than addition, yet it still treats the current reflection as a black-box
unary function.  It leaves unused the repository's generic machinery for
families, W-tree schedules, and transformer-level reflection.

\paragraph{A new unnormalized language.}
Adding impredicative kind universes or unrestricted kind-in-kind would make a
larger syntax easy to state but would invalidate the total evaluator unless a
new normalization theorem were supplied.  In this contest, an unproved
termination claim is not a contender.

\paragraph{Busy Beaver or semantic maximization.}
Maximizing over arbitrary terminating Lean programs, or over machines known to
halt only propositionally, would introduce noncomputability or require an
oracle.  The finite syntactic enumerations in \eqref{eq:fd-reflection} are
essential.

The proposal below instead reuses the certified Fd evaluator and makes a
qualitative jump in the \emph{reflection architecture}: unary reflection is
lifted to family reflection, then to operator reflection, then to a self-height
finite tower of complete operator closures.

%======================================================================
\section{Available infrastructure}

\subsection{Oracle families and their unary code}

BusyLean uses a family
\[
  s:\N\to\N\to\N
\]
as a countable collection of unary oracle functions.  The already defined
\lean{packFamily} maps such a family to one unary oracle by a computable pairing
code.  Hence an Fd program can dynamically select both a family coordinate and
an argument while the metalanguage still passes a single function
$\N\to\N$.

For progress proofs, it is not enough that the object language can in principle
query the family.  The reflector must contain every relevant seed value as an
\emph{explicit} lower bound.  BusyLean's \lean{prefixMax} provides
\[
  \operatorname{prefixMax}(s,n)=\max_{i\le n}s(i,n).
\]
The visibility restriction $i\le n$ is exactly the hypothesis used by the
generic notion of a progressive transformer.

\subsection{Progressive transformers}

A transformer maps oracle families to oracle families:
\[
  T:\Fam\to\Fam.
\]
The existing predicate is
\begin{equation}\label{eq:progressive}
  \Prog(T)\;:\Longleftrightarrow\;
  \forall s,i,o,n,\quad i\le n\Longrightarrow s(i,n)<T(s)(o,n).
\end{equation}
Notice that the output coordinate $o$ is unrestricted.  This asymmetry is
useful: a visible seed coordinate can be placed below any selected output of a
transformed family.

The repository already proves three generic facts.

\begin{enumerate}[leftmargin=*]
\item A progressive transformer remains progressive when iterated down a
      W-tree by \lean{runTransformer}.
\item Its autonomous lift along the universal W-language tree \lean{index}
      remains progressive.
\item \lean{operatorStage} recursively lifts a transformer-valued family along
      another W-tree and preserves progressivity.
\end{enumerate}

These theorems are purely structural.  They do not know whether the primitive
reflection language is System~T, System~F, F$\omega$, kind-polymorphic
F$\omega$, or the new Fd language.

\subsection{Why oracle monotonicity is deliberately absent}

One might hope to prove $s\le s'$ pointwise and infer $T(s)\le T(s')$.  That is
false for the reflected languages: object programs may inspect an oracle value
and branch on it, so increasing an oracle can decrease a program's output.
BusyLean's generic proofs correctly avoid this trap.  Every strict inequality
follows a concrete occurrence of a child or seed in a finite maximum, or a
concrete path through a W-tree.  The new proof follows the same discipline.

%======================================================================
\section{The new family-aware Fd reflection}

\begin{construction}[Family-aware dynamic-kind reflection]\label{con:reflect14}
For a family $s$ and budget $n$, define
\begin{equation}\label{eq:reflect14}
  R_{14}(s,n)=n+1+\max\left(
    \max_{i\le n}s(i,n),\;R^{\Fd}_{\widehat{s}}(n)
  \right),
\end{equation}
where $\widehat{s}=\operatorname{packFamily}(s)$ is its unary code.
\end{construction}

The Lean definition is intentionally direct:
\begin{lstlisting}
def reflect14 (fs : W.Family) (n : Nat) : Nat :=
  n + 1 +
    BigNumber.maxN
      (W.prefixMax (fun i => fs i n) n)
      (Fd.reflectFd (packFamily fs) n)

def fdTransformer : Transformer := fun fs _ n => reflect14 fs n
\end{lstlisting}

All maxima in \eqref{eq:reflect14} are over finite lists, and the call to
$R^{\Fd}$ is total by the incumbent development.  Therefore $R_{14}$ is a
computable function in the ordinary operational sense, not merely a classically
defined natural number.

\begin{lemma}[Visible-family margin]\label{lem:visible}
For every $i\le n$,
\[
  s(i,n)<R_{14}(s,n).
\]
\end{lemma}
\begin{proof}
The value $s(i,n)$ is at most the prefix maximum, which is at most the outer
maximum.  The summand $n+1$ makes the inequality strict.  The Lean theorem is
\lean{oracle\_lt\_reflect14}.
\end{proof}

\begin{lemma}[Fd-program margin]\label{lem:fdprogram}
If $t\in\mathcal{T}_{\Fd}(n)$, then
\[
  \operatorname{value}_{\widehat{s}}(n,t)<R_{14}(s,n).
\]
\end{lemma}
\begin{proof}
The existing theorem puts the program value below
$R^{\Fd}_{\widehat{s}}(n)$; that unary reflection is explicitly below the outer
maximum in \eqref{eq:reflect14}.  The Lean theorem is
\lean{fvalue\_lt\_reflect14}.
\end{proof}

\begin{proposition}[Progressivity]\label{prop:base-progressive}
The transformer \lean{fdTransformer} is progressive.
\end{proposition}
\begin{proof}
The output coordinate is ignored, and \cref{lem:visible} is exactly the
predicate in \eqref{eq:progressive}.  This is
\lean{fdTransformer\_progressive}.
\end{proof}

%======================================================================
\section{A stronger autonomous schedule}

Contender~12 introduced a W-tree schedule obtained from transparent ordinal-like
syntax.  The syntax has constructors for zero, successor, and countable limit,
plus structural operations resembling addition, multiplication, omega powers,
derivatives, a common-critical closure, and higher derivatives.  No quotient by
ordinal equality is required; only the generated tree matters.

Let $T_{12}$ denote the existing \lean{tree12}.  The new fixed schedule is
\begin{equation}\label{eq:tree14}
  T_{14}=\operatorname{realize}
  \left(\operatorname{hyper}_{42}(\operatorname{encode}(T_{12}))\right).
\end{equation}
In Lean:
\begin{lstlisting}
def tree14 : WTree Nat :=
  Autonomous.realize
    (Autonomous.hyper 42 (Autonomous.encode tree12))
\end{lstlisting}

This is a second autonomous 42-level closure placed \emph{above the whole
schedule} used by Contender~12.  The phrase ``above'' here concerns the
syntactic embedding and the explicit child structure.  The formal comparison
with Contender~13 does not depend on assigning a quotient-level ordinal to
$T_{14}$, so no delicate ordinal-identity claim enters the trusted argument.

%======================================================================
\section{Finite powers of complete operator stages}

Fix a W-tree $a$ and a base transformer $B$.  Define transformer powers by
\begin{equation}\label{eq:oppower}
  P^{a,B}_0=B,
  \qquad
  P^{a,B}_{k+1}=\OpStage(P^{a,B}_k,a).
\end{equation}
The executable Lean definition is only four lines:
\begin{lstlisting}
def operatorPower (tree : WTree Nat) (base : Transformer) :
    Nat -> Transformer
  | 0 => base
  | k + 1 => operatorStage (operatorPower tree base k) tree
\end{lstlisting}
This is a finite recursion on $k$.  Even when $k$ is astronomically large, the
function is computable in principle.  The contest does not require feasible
numerical evaluation; it requires a terminating definition and a feasible
\emph{type-checking} proof.

\begin{theorem}[Preservation of progressivity]\label{thm:power-progressive}
If $B$ is progressive, then every $P^{a,B}_k$ is progressive.
\end{theorem}
\begin{proof}
Induct on $k$.  The zero case is the hypothesis.  The successor case is exactly
the repository theorem \lean{operatorStage\_progressive}.  The formal theorem
is \lean{operatorPower\_progressive}.
\end{proof}

The next lemma records a stronger structural margin that is useful for program
domination.

\begin{lemma}[Base below one complete stage]\label{lem:base-below-stage}
Let $B$ be progressive.  For every W-tree $a$, family $s$, and $i\le n$,
\[
  B(s)(i,n)<\OpStage(B,a)(s)(o,n)
\]
for every output coordinate $o$.
\end{lemma}
\begin{proof}
At a leaf, \lean{operatorStage} is the autonomous lift, and the existing theorem
\lean{transformer\_lt\_lift\_at} gives the claim.  At a node, recurse into child
zero and then use \lean{operatorStage\_child\_lt} to enter the parent.  This
follows an explicit path through the tree and assumes no oracle monotonicity.
The Lean theorem is \lean{transformer\_lt\_operatorStage}.
\end{proof}

\begin{corollary}[Strict power chain]\label{cor:power-chain}
For visible input coordinates,
\[
  P^{a,B}_k(s)(i,n)<P^{a,B}_{k+1}(s)(o,n).
\]
Moreover, $B$ lies below every positive power.
\end{corollary}
\begin{proof}
Apply \cref{lem:base-below-stage} first to $P_k$, then induct on the number of
powers.  The formal theorems are \lean{operatorPower\_step\_lt} and
\lean{transformer\_lt\_operatorPower}.
\end{proof}

%======================================================================
\section{The self-height diagonal}

Define the complete incumbent family
\begin{equation}\label{eq:family13}
  G(n)=F_{13}(F_{13}(n)).
\end{equation}
Thus $G(42)$ is definitionally $C_{13}$, but the new source never mentions the
constant \lean{contender\_13}.  Let $\underline{G}$ be the constant-coordinate
family $\underline{G}(i,n)=G(n)$.  Use $B=\lean{fdTransformer}$ and
$a=T_{14}$ in \eqref{eq:oppower}.

\begin{definition}[Contender 14 family]\label{def:f14}
\begin{equation}\label{eq:f14}
  F_{14}(n)
  =P^{T_{14},B}_{G(n)+1}(\underline{G})(G(n),G(n)).
\end{equation}
The proposed closed number is
\begin{equation}\label{eq:c14}
  C_{14}=F_{14}(F_{14}(42)).
\end{equation}
\end{definition}

The three appearances of $G(n)$ in \eqref{eq:f14} have distinct roles:

\begin{enumerate}[leftmargin=*]
\item $G(n)+1$ is the number of \emph{complete operator-stage rounds};
\item the family argument supplies $G$ as the oracle seed at every coordinate;
\item $G(n)$ is both the selected output coordinate and the numerical budget.
\end{enumerate}

The Lean code is correspondingly explicit:
\begin{lstlisting}
def family13 (n : Nat) : Nat := F13 (F13 n)

def F14 (n : Nat) : Nat :=
  operatorPower tree14 fdTransformer (family13 n + 1)
    (fun _ => family13) (family13 n) (family13 n)

def contender_14 : Nat := F14 (F14 42)
\end{lstlisting}

\subsection{The strict comparison}

\begin{lemma}[The incumbent family exceeds its input]\label{lem:g-input}
For all $n$, $n<G(n)$.
\end{lemma}
\begin{proof}
The existing Fd input-margin theorem gives
$n<F_{13}(n)<F_{13}(F_{13}(n))$.  The Lean theorem is
\lean{input\_lt\_family13}.
\end{proof}

\begin{lemma}[Self-height margin]\label{lem:self-height}
For all $n$,
\[
  G(G(n))<F_{14}(n).
\]
\end{lemma}
\begin{proof}
By \cref{thm:power-progressive}, the final transformer in
\eqref{eq:f14} is progressive.  Apply progressivity to seed coordinate zero,
output coordinate $G(n)$, and numerical input $G(n)$.  The seed value is
$\underline{G}(0,G(n))=G(G(n))$.  This is the formal theorem
\lean{family13\_twice\_lt\_F14}.
\end{proof}

\begin{theorem}[Pointwise domination]\label{thm:pointwise}
For every $n$,
\[
  G(n)<F_{14}(n).
\]
\end{theorem}
\begin{proof}
By \cref{lem:g-input} at input $G(n)$,
$G(n)<G(G(n))$.  Compose with \cref{lem:self-height}.  This is
\lean{family13\_lt\_F14}.
\end{proof}

\begin{theorem}[Required contest comparison]\label{thm:main}
\[
  C_{13}<C_{14}.
\]
\end{theorem}
\begin{proof}
By \cref{thm:pointwise} at $42$ and the input margin for $F_{14}$,
\[
  C_{13}=G(42)
  <F_{14}(42)
  <F_{14}(F_{14}(42))=C_{14}.
\]
The Lean theorem is
\lean{contender\_13\_lt\_contender\_14}.
\end{proof}

The comparison is therefore not a fragile numerical estimate.  It is a short
composition of strict inequalities generated by visible occurrences in the
construction.

%======================================================================
\section{Program-level domination}

The construction retains a semantic margin beyond the predecessor constant.
Let $t$ be any accepted dynamic-kind term at budget $G(n)$.  By
\cref{lem:fdprogram},
\[
  \operatorname{value}_{\widehat{\underline{G}}}(G(n),t)
  <R_{14}(\underline{G},G(n)).
\]
By \cref{cor:power-chain}, the base reflection output lies below the positive
operator power in \eqref{eq:f14}.  Hence:

\begin{theorem}[Diagonal Fd-program domination]\label{thm:program}
If $t\in\mathcal{T}_{\Fd}(G(n))$, then
\[
  \operatorname{value}_{\widehat{\underline{G}}}(G(n),t)<F_{14}(n).
\]
\end{theorem}

The Lean theorem is \lean{fvalue\_lt\_F14}.  This matters conceptually.  The
new number does not merely iterate the old numerical family; its first
reflection layer already majorizes the entire bounded Fd programming language
relative to a packed family oracle, and the result is then passed through a
self-height tower of complete operator closures.

%======================================================================
\section{Constructivity and contest compliance}

\begin{longtable}{>{\raggedright\arraybackslash}p{0.23\textwidth}
                  >{\raggedright\arraybackslash}p{0.70\textwidth}}
\toprule
\textbf{Requirement} & \textbf{How the proposal satisfies it}\\
\midrule
\endfirsthead
\toprule
\textbf{Requirement} & \textbf{How the proposal satisfies it}\\
\midrule
\endhead
Closed natural number & \lean{contender\_14} has type \lean{Nat} and no free
parameters.\\
Computable in principle & All recursions are structural: Fd normalization is
already certified; maxima are finite; W-tree traversals use inductive
recursors; \lean{operatorPower} is primitive recursion on a natural number.\\
No previous constant & The definition uses \lean{family13 n := F13 (F13 n)}
and never references \lean{contender\_13} or any earlier closed contender.\\
Qualitative strengthening & A family-aware Fd reflection is lifted to a
transformer, then to an operator stage, then iterated for a previous-family-sized
height along a schedule above \lean{tree12}.\\
Formal strict comparison & \lean{contender\_13\_lt\_contender\_14} follows from
the pointwise theorem \lean{family13\_lt\_F14}.\\
Axiom discipline & The audit file prints the axiom closure of the number and
key theorems and rejects dependencies on earlier contender constants.\\
No forbidden implementation device & The static audit rejects
\lean{noncomputable}, \lean{partial}, \lean{unsafe}, \lean{opaque}, external
implementations, \lean{native\_decide}, \lean{sorry}, custom axioms, and
\lean{Classical.choice}.\\
No giant-term normalization in proofs & Proofs use named margin lemmas and
explicit paths.  They avoid simplification tactics that could unfold a
contender-sized recursion.\\
Time budgets & The supplied clean-build harness separately times the previous
accepted stack, the new definition module, and the new comparison module under
the repository's 15-second and 60-second limits.\\
\bottomrule
\end{longtable}

%======================================================================
\section{Heuristic growth analysis}

This section is explanatory and is \emph{not} part of the formal theorem.
Exact proof-theoretic calibration of the custom Fd language and of the generic
operator construction would require a separate ordinal analysis.

Let $\Lambda_{\Fd}$ denote the effective growth strength of one dynamic-kind
reflection step, and let $\alpha_{12}$ and $\alpha_{14}$ denote heuristic
heights associated with \lean{tree12} and \lean{tree14}.  The source structure
suggests
\[
  \alpha_{14}\approx H_{42}(\alpha_{12}),
\]
where $H_{42}$ stands only for the 42-level autonomous-derivative syntax used
by the implementation.  It is not asserted to be a canonical ordinal function.

A single Fd reflection, as in Contender~13, adds one bounded-language
majorization at a very large diagonal input.  A complete operator stage does
more: it turns a transformer into a hierarchy of child transformers, takes an
autonomous lift along the W-language's universal index at each node, and then
iterates this operation through the external schedule.  If one writes
$\Phi_{\alpha}$ schematically for the resulting closure on growth indices, the
new family has the rough form
\[
  F_{14}(n)\ \text{behaves like}\
  \Phi_{\alpha_{14}}^{\,G(n)+1}
  (\text{Fd-reflection over }G)(G(n)),
\]
where the superscript denotes finite iteration.  Since the iteration count is
itself $G(n)$, this is a self-height closure, not a fixed finite stack.

Three independent effects separate the proposal from earlier contenders:

\begin{enumerate}[leftmargin=*]
\item the primitive language is Fd, stronger than the Fk language amplified by
      Contenders 11--12;
\item the schedule places an additional autonomous closure above
      Contender~12's entire schedule; and
\item the number of complete operator stages is not $42$ or another fixed
      numeral but the full predecessor-family value $G(n)+1$.
\end{enumerate}

The only size claim needed for acceptance remains the exact Lean inequality in
\cref{thm:main}.  The heuristic picture explains why the margin should be
regarded as substantial rather than cosmetic.

%======================================================================
\section{Lean proof engineering}

\subsection{Never normalize a contender}

The largest practical danger is not logical but elaborational.  A seemingly
innocent \lean{simp [F14]} or an aggressive \lean{change} can ask Lean to
compare terms by reducing a recursion whose index is $G(n)$.  That computation
is impossible in practice even though it terminates mathematically.  The proof
therefore unfolds only small outer wrappers and applies general lemmas whose
parameters remain symbolic.

\subsection{Use explicit lower-bound occurrences}

Every margin theorem follows one of four patterns:

\begin{enumerate}[leftmargin=*]
\item membership in a finite maximum;
\item a visible family coordinate in \lean{prefixMax};
\item a child-to-parent path in a W-tree; or
\item transitivity of already strict inequalities.
\end{enumerate}

This makes the proof robust against the nonmonotonic behavior of programs with
oracle access.

\subsection{Separate definition and proof budgets}

\lean{Contender14.lean} contains only the seven executable definitions.  The
more elaborate structural lemmas live in \lean{Comparison14.lean}.  The check
script first builds the accepted Contender~13 stack, then times the new
definition module against 15 seconds and the proof module against 60 seconds.
The source itself remains small: the definition module is under one hundred
lines, and the proof module is under two hundred lines.

%======================================================================
\section{Validation status and reproducibility}

The package was developed against the exact public signatures in the stated
BusyLean snapshot.  A source-level audit was run successfully: it found all
required definitions and theorems, found no forbidden construct, found no
reference to any earlier \lean{contender\_N} constant in the comment-stripped
definition, and verified a 100-column source limit.

\begin{warning}
The execution container used to prepare this archive did not contain the
repository-pinned Lean~4.34.0 toolchain and could not retrieve its large binary
archive through the available restricted network path.  Consequently, this
package does \emph{not} claim that the two new Lean modules were kernel-built in
that container.  The archive includes \lean{scripts/check14.py} and
\lean{test/Audit14.lean} so that compilation, timing, axiom inspection, and
transitive dependency checks can be performed exactly in a BusyLean checkout.
\end{warning}

The distinction is important.  The mathematical argument and Lean source have
been made as explicit and conservative as possible, but only a successful
kernel run upgrades the proposal from a carefully audited formalization to a
verified repository contender.

\subsection{Reproduction commands}

After copying the overlay files into a checkout at the target snapshot or a
compatible successor:
\begin{lstlisting}[language=bash]
cp BusyLean/Contender14.lean /path/to/BusyLean/BusyLean/
cp BusyLean/Comparison14.lean /path/to/BusyLean/BusyLean/
cp test/Audit14.lean /path/to/BusyLean/test/
cp scripts/check14.py /path/to/BusyLean/scripts/
cd /path/to/BusyLean
python3 scripts/check14.py
\end{lstlisting}
The integration patches additionally add the imports to \lean{BusyLean.lean}
and extend the repository-wide audit.

%======================================================================
\section{Further research}

The present proposal maximizes leverage from already proved infrastructure.
Several deeper successors are natural.

\paragraph{A genuine ladder of dynamic kinds.}
Introduce kinds $\mathtt{dyn}_0,\mathtt{dyn}_1,\ldots$, where
$\mathtt{dyn}_{j+1}$ packages semantic operators built from lower levels.  A
budget-$n$ language could expose the first $n$ levels.  The main challenge is a
uniform candidates model and a fuel-bounded checker whose soundness proof does
not force large computations.

\paragraph{Fd-definable schedules.}
The current schedule is built in Lean from the W-language's universal tree.
A stronger object language with safe large elimination could return tree codes
to Lean, allowing the dynamic-kind language itself to generate the reflection
schedule.

\paragraph{Transfinite operator powers.}
Replace the finite \lean{operatorPower} index by a W-tree recursion whose nodes
carry whole progressive transformers.  This must preserve computability and
must expose enough seed occurrences to prove strictness without assuming
monotonicity.

\paragraph{Exact ordinal analysis.}
Develop a conservative proof-theoretic calibration for dynamic-kind
F$\omega$, its family reflector, and the operator-stage closure.  The article's
current size discussion intentionally stops at structural heuristics.

\paragraph{Machine-checked metatheory of the contest rules.}
Formalize the source-closure audit, forbidden-construct policy, and timing
interface as a reusable Lean or external checker with a small verified core.
This would make future contender acceptance more reproducible.

%======================================================================
\appendix
\section{Complete definition module}

\lstinputlisting{../BusyLean/Contender14.lean}

\section{Complete comparison module}

\lstinputlisting{../BusyLean/Comparison14.lean}

\section{Theorem inventory}

\begin{longtable}{p{0.34\textwidth}p{0.60\textwidth}}
\toprule
\textbf{Lean theorem} & \textbf{Role}\\
\midrule
\endfirsthead
\toprule
\textbf{Lean theorem} & \textbf{Role}\\
\midrule
\endhead
\lean{lt\_reflect14} & Generic bound-to-strict-margin lemma.\\
\lean{oracle\_lt\_reflect14} & Every visible family coordinate is dominated.\\
\lean{packed\_reflection\_lt\_reflect14} & The embedded unary Fd reflection is dominated.\\
\lean{fvalue\_lt\_reflect14} & Every accepted Fd program at the current budget is dominated.\\
\lean{fdTransformer\_progressive} & Discharges the generic operator machinery's hypothesis.\\
\lean{operatorPower\_progressive} & Every finite complete-stage power is progressive.\\
\lean{transformer\_lt\_operatorStage} & Explicit path from a base transformer into one complete stage.\\
\lean{operatorPower\_step\_lt} & Consecutive powers form a strict chain.\\
\lean{transformer\_lt\_operatorPower} & The base lies below every positive power.\\
\lean{input\_lt\_family13} & The incumbent family exceeds its input.\\
\lean{family13\_twice\_lt\_F14} & Direct self-height seed margin.\\
\lean{reflect14\_lt\_F14} & Base family reflection lies below the whole power.\\
\lean{family13\_lt\_F14} & Pointwise domination of the incumbent family.\\
\lean{input\_lt\_F14} & New family exceeds its input.\\
\lean{fvalue\_lt\_F14} & Program-level domination at the diagonal budget.\\
\lean{contender\_13\_lt\_contender\_14} & Required contest theorem.\\
\bottomrule
\end{longtable}

\begingroup
\footnotesize
\begin{thebibliography}{99}
\setlength{\itemsep}{0pt}

\bibitem[Aar99]{aaronson1999}
Scott Aaronson.
\newblock \emph{Who Can Name the Bigger Number?}
\newblock Web essay, 1999.
\newblock \url{https://www.scottaaronson.com/writings/bignumbers.html}.

\bibitem[BR26]{busyLean}
Vladimir Reshetnikov and contributors.
\newblock \emph{BusyLean: Name the Biggest Number, in Lean}.
\newblock GitHub repository, 2026; snapshot
\texttt{0cd13c0c2607b4fcb461090a95f356e538c0fc5b}.
\newblock \url{https://github.com/VladimirReshetnikov/BusyLean}.

\bibitem[DMU21]{deMouraUllrich2021}
Leonardo de Moura and Sebastian Ullrich.
\newblock The Lean 4 theorem prover and programming language.
\newblock In \emph{Automated Deduction -- CADE 28}, LNCS 12699,
pages 625--635. Springer, 2021.

\bibitem[GLT89]{girardLafontTaylor1989}
Jean-Yves Girard, Yves Lafont, and Paul Taylor.
\newblock \emph{Proofs and Types}.
\newblock Cambridge University Press, 1989.

\bibitem[Rey74]{reynolds1974}
John C. Reynolds.
\newblock Towards a theory of type structure.
\newblock In \emph{Programming Symposium}, LNCS 19, pages 408--425, 1974.

\bibitem[Rou14]{rouxContest}
Cody Roux and contributors.
\newblock \emph{Name the Biggest Number}.
\newblock GitHub repository, 2014.
\newblock \url{https://github.com/codyroux/name-the-biggest-number}.

\end{thebibliography}
\endgroup

\end{document}
